\begin{textsample}{NEG Dim 4 – 2009\_c5 – Score -1.00 – 2009\_c5/enade\_2009\_dir\_97289d15da.txt}  \label{ex:f4_neg_086}
Capítulo 5 Distribuição dos Conceitos 5.1 Panorama nacional da distribuição dos conceitos A Tabela 5.1 apresenta a quantidade e o percentual de cursos de DIREITO participantes do ENADE/2009 que se situaram em cada uma das faixas de conceitos, distribuídos de acordo com a Região em que se localizam os cursos das instituições de educação superior, mantenedoras dos cursos responsáveis pela oferta.
O Gráfico 5.1 apresenta a mesma informação sob a forma de diagrama de caixa para uma melhor visualização da dispersão dos conceitos obtidos.
No diagrama de caixa, a caixa propriamente dita fica limitada pelos quartis.
A barra dentro da caixa em negrito representa a mediana.
Linhas verticais unem os quartis aos valores extremos.
Se estes valores extremos se configurarem como valores atípicos, estes aparecem como pontos isolados e as linhas terminam no maior ( ou menor ) valor ainda não considerado atípico.
Eventuais valores atípicos aparecem como pequenos círculos com a Categoria Administrativa correspondente.
Valores considerados muito atípicos são representados por estrelas.
Um diagrama de caixa para poucos dados ou com dados muito concentrados pode ter alguns pontos coincidentes.
Por exemplo, se mais de metade dos cursos de uma Região receberem conceito 5, o valor extremo superior, o quartil superior e a mediana coincidirão.
Se o conjunto só contém um dado este aparece denotado como a barra da mediana.
Se o conjunto apresenta todas as notas iguais, a notação é a mesma, somente uma barra horizontal.
Observando -se os dados da Tabela 5.1, percebe -se que, em termos nacionais, predominaram os cursos que receberam conceito 3, correspondendo respectivamente a 36,8 \% do total nacional.
Exatamente 24,1 \% dos cursos receberam conceito 2, 11,4 \% conceito 4, 3,6 \% o conceito máximo e 2,3 \% dos cursos foram classificados com conceito mínimo.
Dos 968 cursos, 212 ficaram sem conceito, correspondendo a 21,9 \% do total.
Todos os cursos com conceito 5 aparecem como valores atípicos.
Considerando -se separadamente as cinco Regiões ( ver também Gráfico 5.1 ) observa -se que, na Região Norte, 61 cursos participaram, entre os quais o valor modal foi o conceito 3 obtido por 23 cursos ( 37,7 \% do total ).
Nenhum curso da Região obteve o conceito mínimo e 16 cursos ( 26,2 \% ) ficaram sem conceito.
Apenas 4 cursos ( 6,6 \% ) obtiveram o conceito máximo, 11 cursos ( 18,0 \% ) obtiveram conceito 2 e outros 7 cursos ( 11,5 \% ), conceito 4.
Como a concentração no conceito 3 foi muito grande e centralizada, cursos com conceitos abaixo e acima deste valor aparecem como valores atípicos no gráfico.
Esta Região foi uma das que apresentou cursos com a melhor distribuição de conceitos.
Na Região Nordeste, 187 cursos participaram, entre os quais o valor modal foi o conceito 3 obtido por 56 cursos ( 29,9 \% do total ).
Somente 8 cursos ( 4,3 \% ) desta Região obtiveram, conceito máximo, bem como 2 ( 1,1 \% ) cursos obtiveram o conceito mínimo e 58 cursos ( 31,0 \% ) ficaram sem conceito.
O conceito 2 foi obtido por 19,8 \% dos cursos e o conceito 4 por 13,9 \%.
Na Região Sudeste, 409 cursos participaram, entre os quais o valor modal foi o conceito 3 obtido por 144 cursos ( 35,2 \% do total regional ).
O conceito máximo foi alcançado por 12 cursos ( 2,9 \% ), 15 cursos ( 3,7 \% ) obtiveram conceito mínimo e 124 cursos ( 30,3 \% ) ficaram com o conceito 2, 46 cursos ( 11,2 \% ) com conceito 4.
Ficaram sem conceito 68 cursos ( 16,6 \% ).
Na Região Sul, 103 dos 207 cursos avaliados, equivalentes a 49,8 \% do total regional, classificaram -se com conceito 3.
O conceito máximo foi alcançado por 6 cursos ( 2,9 \% ), 2 cursos ( 1,0 \% ) obtiveram conceito mínimo e 41 cursos ( 19,8 \% ) ficaram sem conceito.
Esta Região, juntamente com as Regiões Norte e Nordeste, foi a que apresentou cursos com a melhor distribuição de conceitos.
Na Região Centro-Oeste, 104 cursos participaram, sendo que 28,8 \% ( 30 cursos ) obtiveram conceito 2 e outros 28,8 \%, o conceito 3.
O conceito máximo foi alcançado por 5 cursos ( 4,8 \% ), 3 cursos ( 2,9 \% ) obtiveram conceito mínimo e 29 cursos ( 27,9 \% ) ficaram sem conceito. 5.2 Conceitos por Categoria Administrativa e por Região A Tabela 5.2 apresenta a distribuição dos cursos participantes do ENADE/2009 de DIREITO, por Categoria Administrativa, de acordo com os conceitos por eles alcançados, segundo Região.
O Gráfico 5.2 apresenta a mesma informação sob a forma de diagrama de caixa para uma melhor visualização da dispersão dos conceitos obtidos.
Somente eventuais valores atípicos aparecem com o identificador da Região na qual o curso foi habilitado.
De acordo com as informações apresentadas, a grande maioria, 88,3 \%, é de cursos em instituições privadas, 5,0 \% dos cursos eram vinculados a instituições federais, 4,0 \% a instituições estaduais e os restantes 2,7 \% a instituições municipais.
As instituições federais foram as que apresentaram, em termos nacionais, maior proporção de cursos com conceito 5, equivalentes a 33,3 \% do total dessa categoria.
O conceito modal para esta categoria foi 4, com 41,7 \% dos cursos.
Poucos cursos em instituição federal ( 2 cursos correspondendo a 4,2 \% ) receberam o conceito mínimo, e 2 cursos ( 4,2 \% ) ficaram sem conceito.
Entre as instituições estaduais, 35,9 \% receberam conceito 4.
Nenhum curso em instituição estadual recebeu o conceito mínimo, 25,6 \% receberam o conceito máximo e 8 cursos ( 20,5 \% ) ficaram sem conceito.
Entre as instituições municipais, 38,5 \% receberam conceito 3.
Entre as instituições municipais, 2 cursos ( 2,7 \% ) receberam o conceito mínimo, nenhuma recebeu o conceito máximo e 3 cursos ( 11,5 \% ) ficaram sem conceito.
Os cursos nas instituições privadas também tiveram como conceito modal o valor 3 com uma incidência de 38,8 \%.
Somente 9 cursos ( 1,1 \% ) receberam o conceito máximo, 18 ( 2,1 \% ) receberam o conceito mínimo e 199 ( 23,3 \% ) ficaram sem conceito.
Considerando -se a análise por Região observou -se que em todas as Regiões brasileiras, há representação de cursos em todas categorias administrativas analisadas.
Instituições privadas constituem a maioria em todas as Regiões.
Na Região Norte, o conceito modal para cursos em instituições federais foi o conceito 4 ( 36,4 \% ) e nas privadas o conceito 3 ( 42,6 \% ).
Nas instituições federais, o conceito máximo foi obtido por 3 cursos ( 27,3 \% ).
Participaram 2 cursos de instituições estaduais e 1 curso de instituição municipal.
Dos 2 cursos em instituições estaduais, 1 ficou sem conceito e 1 curso obteve o conceito máximo.
O único curso de instituição municipal obteve o conceito 2.
Nesta Região, nenhum curso de nenhuma instituição ( federal, estadual, municipal e privada ) obteve conceito mínimo e 16 cursos ficaram sem conceito dos quais 1 em instituição federal, 1 em estadual e 14 em instituição privada.
Na Região Nordeste, o conceito modal para cursos avaliados em instituições federais e estaduais foi o conceito 4 ( respectivamente 50,0 \% e 31,1 \% ), nas privadas o conceito 3 ( 33,1 \% ).
Dos 2 cursos em instituição municipal na Região 1 recebeu conceito 4 e o outro ficou sem conceito.
Nesta Região 2 cursos em instituição privada ficaram com conceito mínimo.
O conceito máximo foi obtido somente por 4 cursos em instituições federais e 4 em instituições estaduais.
Sem conceito, ficaram 58 cursos, sendo 6 em estadual, 1 em municipal e 51 em privada.
Na Região Sudeste dos 409 participantes a grande maioria estava em instituições privadas ( 94,4 \% ).
Nesta Região, entre os avaliados, o conceito modal para instituições federais e estaduais foi o conceito 4 ( respectivamente 60,0 \% e 50,0 \% ), nas municipais e nas privadas o conceito 3 ( respectivamente 44,4 \% e 36,0 \% ).
Nesta Região, 67 cursos em instituições privadas e 1 em estadual ficaram sem conceito.
Dois cursos em instituições federais, 1 em estadual e 9 em privada alcançaram o conceito máximo.
O mínimo foi \textbf{alocado} também a poucos cursos, 1 em instituição federal, 1 em municipal e 13 em instituições privadas. 4,2 \% ) receberam o conceito mínimo, e 2 cursos ( 4,2 \% ) ficaram sem conceito.
Entre as instituições estaduais, 35,9 \% receberam conceito 4.
Nenhum curso em instituição estadual recebeu o conceito mínimo, 25,6 \% receberam o conceito máximo e 8 cursos ( 20,5 \% ) ficaram sem conceito.
Entre as instituições municipais, 38,5 \% receberam conceito 3.
Entre as instituições municipais, 2 cursos ( 2,7 \% ) receberam o conceito mínimo, nenhuma recebeu o conceito máximo e 3 cursos ( 11,5 \% ) ficaram sem conceito.
Os cursos nas instituições privadas também tiveram como conceito modal o valor 3 com uma incidência de 38,8 \%.
Somente 9 cursos ( 1,1 \% ) receberam o conceito máximo, 18 ( 2,1 \% ) receberam o conceito mínimo e 199 ( 23,3 \% ) ficaram sem conceito.
Considerando -se a análise por Região observou -se que em todas as Regiões brasileiras, há representação de cursos em todas categorias administrativas analisadas.
Instituições privadas constituem a maioria em todas as Regiões.
Na Região Norte, o conceito modal para cursos em instituições federais foi o conceito 4 ( 36,4 \% ) e nas privadas o conceito 3 ( 42,6 \% ).
Nas instituições federais, o conceito máximo foi obtido por 3 cursos ( 27,3 \% ).
Participaram 2 cursos de instituições estaduais e 1 curso de instituição municipal.
Dos 2 cursos em instituições estaduais, 1 ficou sem conceito e 1 curso obteve o conceito máximo.
O único curso de instituição municipal obteve o conceito 2.
Nesta Região, nenhum curso de nenhuma instituição ( federal, estadual, municipal e privada ) obteve conceito mínimo e 16 cursos ficaram sem conceito dos quais 1 em instituição federal, 1 em estadual e 14 em instituição privada.
Na Região Nordeste, o conceito modal para cursos avaliados em instituições federais e estaduais foi o conceito 4 ( respectivamente 50,0 \% e 31,1 \% ), nas privadas o conceito 3 ( 33,1 \% ).
Dos 2 cursos em instituição municipal na Região 1 recebeu conceito 4 e o outro ficou sem conceito.
Nesta Região 2 cursos em instituição privada ficaram com conceito mínimo.
O conceito máximo foi obtido somente por 4 cursos em instituições federais e 4 em instituições estaduais.
Sem conceito, ficaram 58 cursos, sendo 6 em estadual, 1 em municipal e 51 em privada.
Na Região Sudeste dos 409 participantes a grande maioria estava em instituições privadas ( 94,4 \% ).
Nesta Região, entre os avaliados, o conceito modal para instituições federais e estaduais foi o conceito 4 ( respectivamente 60,0 \% e 50,0 \% ), nas municipais e nas privadas o conceito 3 ( respectivamente 44,4 \% e 36,0 \% ).
Nesta Região, 67 cursos em instituições privadas e 1 em estadual ficaram sem conceito.
Dois cursos em instituições federais, 1 em estadual e 9 em privada alcançaram o conceito máximo.
O mínimo foi \textbf{alocado} também a poucos cursos, 1 em instituição federal, 1 em municipal e 13 em instituições privadas.
Na Região Sul, dos 207 cursos 41 ficaram sem conceito : 1 em municipal e 40 em instituição privada.
Entre os avaliados o conceito 5 foi o modal nas instituições federais, correspondendo a 50 \% do total da Região para esta categoria.
O conceito modal para cursos em instituições estaduais foram os conceitos 4 e 5 ( 50,0 \% cada ), nas municipais e nas privadas o conceito 3 ( respectivamente 60,0 \% e 52,4 \% ).
Dos 6 cursos que alcançaram o conceito máximo, 3 estavam em instituições federais e 3 em estaduais.
O mínimo foi \textbf{alocado} a poucos cursos, 1 em instituição federal e 1 em instituição privada.
Na Região Centro-Oeste, dos 104 cursos 29 ficaram sem conceito : 1 em instituição federal, 1 em municipal e 27 em instituição privada.
Entre os avaliados o conceito modal para cursos em instituições federais foi o conceito 5 ( 44,4 \% ), nas estaduais foram os conceitos 3 e 4 ( 40,0 \% cada ), nas municipais e nas privadas foi o conceito 2 ( 50,0 \% e 32,6 \%, respectivamente ).
Nenhum curso nem em uma instituição privada nem em uma instituição municipal alcançou o conceito máximo nesta Região e 1 curso em uma instituição municipal e 2 em privada obtiveram o conceito mínimo. 5.3 Conceitos por Organização Acadêmica e por Região Na Tabela 5.3 encontra -se a distribuição dos conceitos atribuídos aos cursos participantes do ENADE/2009 de DIREITO, por Organização Acadêmica, segundo as Regiões.
O Gráfico 5.3 apresenta a mesma informação sob a forma de diagrama de caixa para uma melhor visualização da dispersão dos conceitos obtidos.
Somente eventuais valores atípicos apareceram com o identificador da Categoria Administrativa.
Os 968 cursos desta área estão principalmente \textbf{alocados} em Faculdades ( 56,3 \% ), depois em Universidades ( 32,9 \% ) e Centros Universitários ( 10,8 \% ).
Não existem cursos nem em Centros Federais de Educação Tecnológica ( CEFET ) e nem em Institutos Federais de Educação, Ciência e Tecnologia ( IFECT ).
Em todos os tipos de Organização Acadêmica que oferecem cursos em DIREITO que foram avaliados, Universidades, Centros Universitários e Faculdades, o conceito modal foi 3 ( respectivamente 37,4 \%, 47,6 \% e 34,3 \% ).
Proporcionalmente poucos cursos receberam conceito mínimo, 8 em Universidades, 4 em Centros Universitários e 4 em Faculdades ( respectivamente 2,5 \%, 3,8 \% e 1,8 \% ).
Uma proporção também pequena atingiu o conceito máximo : 9,7 \% em Universidades e 0,7 \% em Faculdades.
Nos 3 tipos de Organização Acadêmica existiram 212 cursos que ficaram sem conceito, 22 em Universidade, 4 em Centros Universitários e 186 em Faculdade.
Considerando -se separadamente as cinco Regiões brasileiras, observa -se que, na Região Norte dos 61 cursos, 16 ficaram sem conceito : 2 em Universidades e 14 em Faculdades.
Entre os avaliados, a categoria modal para Universidades foi 4 ( 33,3 \% ), para Centro Universitários foi 2 e 3 ( 44,4 \% cada ) e para Faculdades foi 3 ( 40,5 \% ).
Somente cursos em Universidade nesta Região obtiveram conceito máximo, em número de 4.
O conceito mínimo não foi \textbf{alocado} nesta Região, Na Região Nordeste dos 187 cursos na área de DIREITO, 58 ficaram sem conceito : 7 em Universidades e 51 em Faculdades.
Semelhantemente ao observado na Região Norte, a categoria modal para Universidades foi 4 ( 33,3 \% ), para Centro Universitários foi 2 e 3 ( 50,0 \% cada ) e para Faculdades foi 3 ( 31,2 \% ).
Somente cursos em Universidades nesta Região obtiveram conceito máximo, em número de 8.
Dois cursos em Faculdade obtiveram o conceito mínimo.
Na Região Sudeste, dos 409 cursos na área de DIREITO, 68 ficaram sem conceito : 7 em Universidade, 2 em Centro Universitário, 59 em Faculdade.
Entre os cursos avaliados, a categoria modal nas Universidades foi o conceito 2 ( 35,8 \% ), nos Centros Universitários e Faculdades foi o conceito 3 ( 44,3 \% e 35,0 \%, respectivamente ).
O conceito máximo foi alcançado por poucos cursos e somente em Universidades e Faculdades ( respectivamente 9 e 3 cursos ).
O conceito mínimo foi \textbf{alocado} em 6 cursos em Universidades, 3 em Centros Universitários e 6 cursos em Faculdades.
Na Região Sul dos 207 cursos na área de DIREITO, 41 ficaram sem conceito : 5 em Universidades, 1 em Centro Universitário e 35 em Faculdades.
O conceito modal foi 3 para todos os tipos de Organização Acadêmica : 53,5 \% para Universidades, 63,2 \% para Centros Universitários e 42,7 \% para Faculdades.
O conceito máximo foi alcançado por 5 cursos em Universidades e 1 curso em Faculdade.
O conceito mínimo foi \textbf{alocado} a 2 cursos em Universidades.
Na Região Centro-Oeste, dos 104 cursos na área de DIREITO, 29 ficaram sem conceito : 1 em Universidade, 1 em Centro Universitário e 27 em Faculdades.
A categoria modal para Universidades foi 3 ( 36,0 \% ), para os Centros Universitários foram os conceitos 2 e 3 ( 41,7 \% cada ) e para Faculdades foi o conceito 2 ( 31,3 \% ).
O conceito máximo foi alcançado por 5 cursos em Universidades ( 20,0 \% ).
O conceito mínimo foi \textbf{alocado} a 1 curso em Centro Universitários e 2 cursos em Faculdades.

% matched lemmas: alocar
\end{textsample}
